\section{Finite cylinder vacua}
\label{sec:reference}

The marked contraction uses a finite normalized half-cylinder state.
We specify that state and the local identities supplied by the companion
\cite{Electric}; the next section applies them to select a single polygon.
The magnetic mean and all estimates for its contour integrals are proved
after the finite identity has been established.

\subsection{The finite diagram model}

Slots are vacant or occupied. A diagram pairs some of its occupied
boundary slots by noncrossing curves. Composition glues matching slots;
an occupancy mismatch or a closed loop gives zero. Tensor product
juxtaposes diagrams, and transpose exchanges inputs and outputs.
Set $\lambda=\pi/8$, $\beta=2\lambda$, $q=e^{i\beta}$,
$c=\cos(3\lambda)$, $d=\cos(2\lambda)$ and $C_0=\cos\lambda$.
The two-slot operator $R(v)$ has the following coefficients:
\[
\begin{aligned}
 L(v)&=\sin(2\lambda+v)\sin(3\lambda+v),\\
 A(v)&=d\sin(3\lambda-v)/L(v),&
 U(v)&=d\sin v/L(v),\\
 B(v)&=\sin v\sin(3\lambda-v)/L(v),&
 E(v)&=\sin(3\lambda-v)\sin(2\lambda-v)/L(v),\\
 F(v)&=\sin(v-\lambda)\sin v/L(v).
\end{aligned}
\]
They multiply, respectively, a single vertical connection, a cap or
cup, a single diagonal connection, two vertical connections, and a
cap together with a cup. The empty diagram has coefficient one; each
single vertical or diagonal connection has two positions. The operator
$S_2$ splits one slot into two: a vacancy becomes a vacancy plus
$\rho_{\rm v}$ times a cup, and an occupied strand can take either
output with coefficient $\rho_{\rm v}$. The zero-to-two-slot operator
$S_3$ is a vacancy plus a cup. These conventions also define their
spin representations in Section~\ref{sec:06-magnetic-algebra}.

For a cyclic list $(u_i)_{1\le i\le N}$ write $x_i=e^{2iu_i}$ with
$\sqrt{x_i}=e^{iu_i}$, and set $D_0(x,y)=x^2+y^2+2dxy$.
The normalized vacuum $U_N$ is the vector of half-cylinder matching
weights with empty component one. Its matching basis has disk
connectivity: compactify the vacant open end to a point and allow no
strand to end there. This retains the ordinary noncrossing partial
matchings even when a subsequent gluing produces a winding polygon.
At physical parameters
$u_i\in\{\lambda,2\lambda\}$ its series with vacant far boundary
converges absolutely. Algebraically, $U_N=V_N/h_N$, where $V_N$ is
homogeneous of degree $N(N-1)/2$ and
\begin{equation}\label{eq:vacuum-pfaffian}
 h_N=\prod_{i<j}\frac{D_0(x_i,x_j)}{x_i-x_j}
       \operatorname{Pf}\left[\frac{x_i^2-x_j^2}{D_0(x_i,x_j)}\right].
\end{equation}
For odd $N$ append a last column of ones and its negative transpose;
$h_0=h_1=1$. In each slot, the vacant part of $V_N$ has degree at
most $N-1$, and the occupied part is $\sqrt{x_i}$ times a polynomial
of degree at most $N-2$.

The finite identities we use are the Yang--Baxter relation, unitarity,
$R(j\lambda)=S_jS_j^{\mathsf t}$ for $j=2,3$, and the corresponding
vacuum braid and splitting identities
\[
\begin{aligned}
 U_N(\ldots,u_i,u_{i+1},\ldots)
 &=R_i(u_{i+1}-u_i)U_N(\ldots,u_{i+1},u_i,\ldots),\\
 U_N\big|_{u_{i+1}=u_i+j\lambda}&=S_jU_{\rm red}.
\end{aligned}
\]
The merge replaces $(q^{-1}x,qx)$ by $x$; the deletion removes
$(x,q^3x)$. The vacuum is cyclically covariant and covariant under
reversing the row while replacing $u_i$ by $3\lambda-u_i$.
Finally
\[
 h_N(x_1^{-1},\ldots,x_N^{-1})
 =h_N(x_1,\ldots,x_N)\prod_i x_i^{-(N-1)}.
\]
These local, homogeneous-vacuum, and stationary-cylinder statements
are supplied by \cite[Proposition~2.2, Lemma~2.3, and
Proposition~2.6]{Electric}. They hold as rational identities
at generic complex parameters and continue at regular coincidences.
In particular the right-half state is the same vacuum on the list
$(3\lambda-u_i)_i$. The distinction between the polynomial $h_N$
and the normalized physical state $U_N$ will matter in every marked
identity below.

We record the scalar factors of the two reductions, since the marked
contraction will be cleared on its original site list while the
splitting identities act on longer lists. Define the centered scalar
\[
 \widetilde h_j(x_1,\ldots,x_j)
 =h_j(x_1,\ldots,x_j)\prod_{i=1}^j x_i^{-(j-1)/2}.
\]
For a spectator list $Z$ of length $j-2$, the scalar reductions in
\cite[Proposition~2.2 and the proof of Lemma~2.3]{Electric} give
\begin{align}
 \frac{\widetilde h_j(Z,q^{-1}x,qx)}
      {\widetilde h_{j-1}(Z,x)}
 &=2d\prod_{z\in Z}\frac{z+x}{\sqrt{zx}},
 \label{eq:vacuum-centered-merge}\\
 \frac{\widetilde h_j(Z,x,y)}{\widetilde h_{j-2}(Z)}
 &=2c\prod_{z\in Z}
       \frac{(z-q^{-3}x)(z-q^6x)}{z\sqrt{xy}},
 \qquad y=q^3x.
 \label{eq:vacuum-centered-deletion}
\end{align}
Here and below the square roots follow the fixed spectral arguments.
Indeed the uncentered factors are respectively
$(q^{-1}+q)x\prod_{z\in Z}(z+x)$ and
$(x+y)\prod_{z\in Z}(z-q^{-3}x)(z-q^6x)$; centering and
$x+y=2c\sqrt{xy}$ give the displayed formulas. They hold first
where the scalar quotients are defined and then as cleared identities.
